% GENERATED FILE. Edit canonical lessons, then run npm run generate:books.
% canonical-source-hash: fnv1a64:13c1a30a6b3e59a9
% canonical-lessons: MR-C391-ghatne, MR-C391-tikne, MR-C391-takka, MR-C391-praman, MR-C391-star

\chapter{Rising and Falling}
\label{ch:mr-rising-and-falling}
\hlchaptermodality{391}

\begin{chapteropening}
\textbf{By the end of this chapter:} \emph{I can say what went up, what went down and by how much.}

Complete the last lesson of chapter 391: I can say what went up, what went down and by how much.
\end{chapteropening}

\section[\texorpdfstring{\mr{घटणे} (ghaṭṇe)}{घटणे (ghaṭṇe)}]{\mr{घटणे} (ghaṭṇe) — to decline; to decrease}

\label{lesson:MR-C391-ghatne}



\begin{quote}
Before the new one: say the Marathi for particular, then the Marathi for apart from.
\end{quote}

\subsection*{You'll want to know: \mr{घटणे}}
\textbf{\mr{घटणे}} — \emph{ghaṭṇe} — \textquotedblleft{}to decline; to decrease\textquotedblright{}.

\begin{quote}
\mr{थंडीत} \mr{भाज्यांचे} \mr{दर} \mr{घटले}.
\end{quote}

\begin{quote}
\emph{thaṇḍīt bhājyāñce dar ghaṭle.}
\end{quote}

\begin{quote}
Vegetable prices fell in the cold season.
\end{quote}

Past (he …): \mr{तो} \mr{घटला}, \emph{to ghaṭlā}.

Intransitive: what decreases is the subject.

\begin{cousinweb}
\textbf{\mr{घटणे}} (\emph{ghaṭṇe}) is inherited: CDIAL lists Marathi \emph{ghaṭṇẽ}, \textquotedblleft{}to decrease\textquotedblright{}, under \emph{*ghaṭṭati}, \textquotedblleft{}decreases, is wanting\textquotedblright{} (CDIAL 4415).
\end{cousinweb}

\begin{grammarlens}[title={What you've built}]
One more word for saying what happens if.
\end{grammarlens}

\subsection*{Your turn}
\begin{itemize}
\raggedright
  \item \emph{Say it:} \emph{ghaṭṇe}
  \item \emph{Say it:} \emph{ghaṭṇe}, once more
  \item \emph{Say it:} \emph{thaṇḍīt bhājyāñce dar ghaṭle.} — Vegetable prices fell in the cold season.
  \item \emph{Recall:} say the Marathi for particular, then the Marathi for apart from, then say \emph{ghaṭṇe} again
\end{itemize}

\begin{tcolorbox}[breakable,colback=teal!4,colframe=teal!35!black,title={Before you move on}]
What does \mr{घटणे} mean? (\textquotedblleft{}To decline; to decrease\textquotedblright{}.) Say it once more.
\end{tcolorbox}
\section[\texorpdfstring{\mr{टिकणे} (ṭikṇe)}{टिकणे (ṭikṇe)}]{\mr{टिकणे} (ṭikṇe) — to hold up; to last; to survive}

\label{lesson:MR-C391-tikne}



\begin{quote}
Before the new one: say the Marathi for apart from, then the Marathi for to decline.
\end{quote}

\subsection*{You'll want to know: \mr{टिकणे}}
\textbf{\mr{टिकणे}} — \emph{ṭikṇe} — \textquotedblleft{}to hold up; to last; to survive\textquotedblright{}.

\begin{quote}
\mr{ही} \mr{चप्पल} \mr{दोन} \mr{वर्षं} \mr{टिकली}.
\end{quote}

\begin{quote}
\emph{hī cappal don varṣaṁ ṭiklī.}
\end{quote}

\begin{quote}
These sandals lasted two years.
\end{quote}

Past (he …): \mr{तो} \mr{टिकला}, \emph{to ṭiklā}.

\mr{टिकाऊ}, durable.

\begin{cousinweb}
\textbf{\mr{टिकणे}} (\emph{ṭikṇe}) is inherited: CDIAL lists Marathi \emph{ṭikṇẽ}, \textquotedblleft{}to stay\textquotedblright{}, under \emph{*ṭakk-} / \emph{*ṭikk-}, \textquotedblleft{}remain, stop\textquotedblright{} (CDIAL 5420).
\end{cousinweb}

\begin{grammarlens}[title={What you've built}]
One more word for saying what happens if.
\end{grammarlens}

\subsection*{Your turn}
\begin{itemize}
\raggedright
  \item \emph{Say it:} \emph{ṭikṇe}
  \item \emph{Say it:} \emph{ṭikṇe}, once more
  \item \emph{Say it:} \emph{hī cappal don varṣaṁ ṭiklī.} — These sandals lasted two years.
  \item \emph{Recall:} say the Marathi for apart from, then the Marathi for to decline, then say \emph{ṭikṇe} again
\end{itemize}

\begin{tcolorbox}[breakable,colback=teal!4,colframe=teal!35!black,title={Before you move on}]
What does \mr{टिकणे} mean? (\textquotedblleft{}To hold up; to last; to survive\textquotedblright{}.) Say it once more.
\end{tcolorbox}
\section[\texorpdfstring{\mr{टक्का} (ṭakkā)}{टक्का (ṭakkā)}]{\mr{टक्का} (ṭakkā) — a percent}

\label{lesson:MR-C391-takka}



\begin{quote}
Before the new one: say the Marathi for to decline, then the Marathi for to hold up.
\end{quote}

\subsection*{You'll want to know: \mr{टक्का}}
\textbf{\mr{टक्का}} — \emph{ṭakkā} — \textquotedblleft{}a percent\textquotedblright{}.

\begin{quote}
\mr{या} \mr{वर्षी} \mr{परीक्षेत} \mr{नव्वद} \mr{टक्के} \mr{मुलं} \mr{उत्तीर्ण} \mr{झाली}.
\end{quote}

\begin{quote}
\emph{yā varṣī parīkṣet navvad ṭakke mulaṁ uttīrṇa jhālī.}
\end{quote}

\begin{quote}
This year ninety percent of the children passed the exam.
\end{quote}

Masculine. Plural: \mr{टक्के}, \emph{ṭakke}.

Plural \mr{टक्के}: \mr{पाच} \mr{टक्के}, five percent.

\begin{cousinweb}
\textbf{\mr{टक्का}} (\emph{ṭakkā}) is usually traced to Sanskrit \emph{ṭaṅka}, \textquotedblleft{}a stamped coin\textquotedblright{}, under which CDIAL lists Marathi \emph{ṭakā}, \textquotedblleft{}a two-paisa coin\textquotedblright{} (CDIAL 5426); the sense one in a hundred came later, and Gujarati \emph{ṭako} has it too.
\end{cousinweb}

\begin{grammarlens}[title={What you've built}]
One more word for saying what happens if.
\end{grammarlens}

\subsection*{Your turn}
\begin{itemize}
\raggedright
  \item \emph{Say it:} \emph{ṭakkā}
  \item \emph{Say it:} \emph{ṭakkā}, once more
  \item \emph{Say it:} \emph{yā varṣī parīkṣet navvad ṭakke mulaṁ uttīrṇa jhālī.} — This year ninety percent of the children passed the exam.
  \item \emph{Recall:} say the Marathi for to decline, then the Marathi for to hold up, then say \emph{ṭakkā} again
\end{itemize}

\begin{tcolorbox}[breakable,colback=teal!4,colframe=teal!35!black,title={Before you move on}]
What does \mr{टक्का} mean? (\textquotedblleft{}A percent\textquotedblright{}.) Say it once more.
\end{tcolorbox}
\section[\texorpdfstring{\mr{प्रमाण} (pramāṇ)}{प्रमाण (pramāṇ)}]{\mr{प्रमाण} (pramāṇ) — a proportion; a rate}

\label{lesson:MR-C391-praman}



\begin{quote}
Before the new one: say the Marathi for to hold up, then the Marathi for a percent.
\end{quote}

\subsection*{You'll want to know: \mr{प्रमाण}}
\textbf{\mr{प्रमाण}} — \emph{pramāṇ} — \textquotedblleft{}a proportion; a rate\textquotedblright{}.

\begin{quote}
\mr{शहरात} \mr{अपघातांचं} \mr{प्रमाण} \mr{वाढलं} \mr{आहे}.
\end{quote}

\begin{quote}
\emph{śahrāt apghātāñcaṁ pramāṇ vāḍhlaṁ āhe.}
\end{quote}

\begin{quote}
The rate of accidents in the city has gone up.
\end{quote}

Neuter.

Also a standard: \mr{प्रमाणभाषा}, the standard language.

\begin{cousinweb}
\textbf{\mr{प्रमाण}} (\emph{pramāṇ}) is a tatsama, taken from Sanskrit \emph{pramāṇa}, \textquotedblleft{}measure, standard, proof\textquotedblright{} (CDIAL 8723).
\end{cousinweb}

\begin{grammarlens}[title={What you've built}]
One more word for saying what happens if.
\end{grammarlens}

\subsection*{Your turn}
\begin{itemize}
\raggedright
  \item \emph{Say it:} \emph{pramāṇ}
  \item \emph{Say it:} \emph{pramāṇ}, once more
  \item \emph{Say it:} \emph{śahrāt apghātāñcaṁ pramāṇ vāḍhlaṁ āhe.} — The rate of accidents in the city has gone up.
  \item \emph{Recall:} say the Marathi for to hold up, then the Marathi for a percent, then say \emph{pramāṇ} again
\end{itemize}

\begin{tcolorbox}[breakable,colback=teal!4,colframe=teal!35!black,title={Before you move on}]
What does \mr{प्रमाण} mean? (\textquotedblleft{}A proportion; a rate\textquotedblright{}.) Say it once more.
\end{tcolorbox}
\section[\texorpdfstring{\mr{स्तर} (star)}{स्तर (star)}]{\mr{स्तर} (star) — a level}

\label{lesson:MR-C391-star}



\begin{quote}
Before the new one: say the Marathi for a percent, then the Marathi for a proportion.
\end{quote}

\subsection*{You'll want to know: \mr{स्तर}}
\textbf{\mr{स्तर}} — \emph{star} — \textquotedblleft{}a level\textquotedblright{}.

\begin{quote}
\mr{नदीतील} \mr{पाण्याचा} \mr{स्तर} \mr{वाढला}.
\end{quote}

\begin{quote}
\emph{nadītīl pāṇyācā star vāḍhlā.}
\end{quote}

\begin{quote}
The water level in the river rose.
\end{quote}

Masculine.

Also a layer, a tier: \mr{समाजाचे} \mr{स्तर}.

\begin{cousinweb}
\textbf{\mr{स्तर}} (\emph{star}) is a tatsama, taken from Sanskrit \emph{stara}, \textquotedblleft{}layer\textquotedblright{} (CDIAL 13685); the same entry lists an inherited Marathi \emph{thar}, \textquotedblleft{}layer\textquotedblright{}.
\end{cousinweb}

\begin{grammarlens}[title={What you've built}]
One more word for saying what happens if.
\end{grammarlens}

\subsection*{Your turn}
\begin{itemize}
\raggedright
  \item \emph{Say it:} \emph{star}
  \item \emph{Say it:} \emph{star}, once more
  \item \emph{Say it:} \emph{nadītīl pāṇyācā star vāḍhlā.} — The water level in the river rose.
  \item \emph{Recall:} say the Marathi for a percent, then the Marathi for a proportion, then say \emph{star} again
\end{itemize}

\begin{tcolorbox}[breakable,colback=teal!4,colframe=teal!35!black,title={Before you move on}]
What does \mr{स्तर} mean? (\textquotedblleft{}A level\textquotedblright{}.) Say it once more.
\end{tcolorbox}
